% Options for packages loaded elsewhere
\PassOptionsToPackage{unicode}{hyperref}
\PassOptionsToPackage{hyphens}{url}
\documentclass[
]{article}
\usepackage{xcolor}
\usepackage[margin=1in]{geometry}
\usepackage{amsmath,amssymb}
\setcounter{secnumdepth}{-\maxdimen} % remove section numbering
\usepackage{iftex}
\ifPDFTeX
  \usepackage[T1]{fontenc}
  \usepackage[utf8]{inputenc}
  \usepackage{textcomp} % provide euro and other symbols
\else % if luatex or xetex
  \usepackage{unicode-math} % this also loads fontspec
  \defaultfontfeatures{Scale=MatchLowercase}
  \defaultfontfeatures[\rmfamily]{Ligatures=TeX,Scale=1}
\fi
\usepackage{lmodern}
\ifPDFTeX\else
  % xetex/luatex font selection
\fi
% Use upquote if available, for straight quotes in verbatim environments
\IfFileExists{upquote.sty}{\usepackage{upquote}}{}
\IfFileExists{microtype.sty}{% use microtype if available
  \usepackage[]{microtype}
  \UseMicrotypeSet[protrusion]{basicmath} % disable protrusion for tt fonts
}{}
\makeatletter
\@ifundefined{KOMAClassName}{% if non-KOMA class
  \IfFileExists{parskip.sty}{%
    \usepackage{parskip}
  }{% else
    \setlength{\parindent}{0pt}
    \setlength{\parskip}{6pt plus 2pt minus 1pt}}
}{% if KOMA class
  \KOMAoptions{parskip=half}}
\makeatother
\usepackage{graphicx}
\makeatletter
\newsavebox\pandoc@box
\newcommand*\pandocbounded[1]{% scales image to fit in text height/width
  \sbox\pandoc@box{#1}%
  \Gscale@div\@tempa{\textheight}{\dimexpr\ht\pandoc@box+\dp\pandoc@box\relax}%
  \Gscale@div\@tempb{\linewidth}{\wd\pandoc@box}%
  \ifdim\@tempb\p@<\@tempa\p@\let\@tempa\@tempb\fi% select the smaller of both
  \ifdim\@tempa\p@<\p@\scalebox{\@tempa}{\usebox\pandoc@box}%
  \else\usebox{\pandoc@box}%
  \fi%
}
% Set default figure placement to htbp
\def\fps@figure{htbp}
\makeatother
\setlength{\emergencystretch}{3em} % prevent overfull lines
\providecommand{\tightlist}{%
  \setlength{\itemsep}{0pt}\setlength{\parskip}{0pt}}
\usepackage{booktabs}
\usepackage{longtable}
\usepackage{array}
\usepackage{multirow}
\usepackage{wrapfig}
\usepackage{float}
\usepackage{colortbl}
\usepackage{pdflscape}
\usepackage{tabularx}
\usepackage{xltabular}
\usepackage{threeparttable}
\usepackage{threeparttablex}
\usepackage[normalem]{ulem}
\usepackage{makecell}
\usepackage{xcolor}
\usepackage{bookmark}
\IfFileExists{xurl.sty}{\usepackage{xurl}}{} % add URL line breaks if available
\urlstyle{same}
\hypersetup{
  pdftitle={ANPSA Conference 2026 - Notes},
  pdfauthor={Peter GS},
  hidelinks,
  pdfcreator={LaTeX via pandoc}}

\title{ANPSA Conference 2026 - Notes}
\author{Peter GS}
\date{2026-08-26}

\begin{document}
\maketitle

\section{Desert ecology - Steve
Morton}\label{desert-ecology---steve-morton}

Major forces at play in shaping arid Australia:

\begin{enumerate}
\def\labelenumi{\arabic{enumi}.}
\item
  Time

  \begin{itemize}
  \tightlist
  \item
    continental drift into mid-latitudes, where arid conditions exist
  \item
    no geological rejuvenation over this time (e.g.~volcanic activity ?)
  \item
    landscape has been formed by cycles of erosion and deposition
  \end{itemize}
\item
  Climate

  \begin{itemize}
  \tightlist
  \item
    extreme variability is the norm
  \item
    east arid zone

    \begin{itemize}
    \tightlist
    \item
      firmer red earths soil type, which resists complete infiltration
      of rainfall and allows surface runoff (Note he said elsewhere that
      red earths allow water to infiltrate deeply. How does this square
      up?)
    \end{itemize}
  \item
    west arid zone

    \begin{itemize}
    \tightlist
    \item
      soil has much higher sand content and water completely infiltrates
      into it
    \end{itemize}
  \end{itemize}
\end{enumerate}

Ephemeral plants use water that's available in the short term.

Perennial plants use water that they can draw on over the long term.

Mulga communities sum up the plant ecology of arid Australia - unclear
to me. (This is the idea in his book about bands of mulga passing on
water and nutrients in some way that gives them a double dose.)

Perennial plants photosynthesise a lot to produce carbohydrate which
they put into root growth so theycan tap into water deep in the soil.

As a lot of plant matter is produced, fire hazard increases.

\{Check out the ecological flowchart - probabky in his book.\}

Ants are everywhere and play a big role in plant propagation (moving
seed around) and plant protection (preying on insects that damage
plants).

What are the main factors that give inland plants their characteristics?

\begin{itemize}
\tightlist
\item
  little rain
\item
  occasional large falls of rain
\item
  infertile soils
\end{itemize}

\emph{Outstanding questions}

\begin{itemize}
\tightlist
\item
  Root structure in soil
\item
  Fertility of desert soils

  \begin{itemize}
  \tightlist
  \item
    average figures are probably not useful as local variability is so
    high
  \end{itemize}
\item
  Role of fungi, nutrient cycling and organisms that consume plants
\item
  Transfer of carbohydrate by ants

  \begin{itemize}
  \tightlist
  \item
    there are few, if any, measures of the importance of ants in moving
    carbohydrate around
  \end{itemize}
\end{itemize}

\section{Biogeography of Central Australian
flora}\label{biogeography-of-central-australian-flora}

Overview

\begin{itemize}
\tightlist
\item
  species groups
\item
  notable absences from these species groups
\item
  origins
\end{itemize}

Characteristics of Central Australia

\begin{itemize}
\tightlist
\item
  varied topography
\item
  large desert areas
\item
  several different bioregions
\end{itemize}

Arid southern Northern Territory has a very large number of species
despite the harsh climate (i.e.~low rainfall, low nutrient soil)

Simpson Desert has about 100 species, with maybe only one or two endemic
species

\begin{itemize}
\tightlist
\item
  this may be due to the Simpson Desert being a relatively young desert
\end{itemize}

Types of desert environment

\begin{itemize}
\item
  sand dunes

  \begin{itemize}
  \tightlist
  \item
    these occur mostly in deserts
  \end{itemize}
\item
  alluvial clay plains

  \begin{itemize}
  \tightlist
  \item
    these provide varied habitats
  \item
    most important are the Burt Plains
  \end{itemize}
\item
  gibber plains

  \begin{itemize}
  \tightlist
  \item
    not common in the Norther Territory - there are some in the Finke
    River bioregion
  \item
    they provide a variety of habitats, which vary with salt level and
    soil
  \end{itemize}
\item
  saltpans and lakes

  \begin{itemize}
  \tightlist
  \item
    for example, Lake Amadeus
  \item
    in saline environments, there are a range of plants adapted to these
    conditions, including some chenopods, daisies, peas and samphires
  \end{itemize}
\item
  rivers and ephemeral swamps

  \begin{itemize}
  \tightlist
  \item
    for example, Hale R., Plenty R., Finke R., the latter being the
    oldest river system in the world
  \item
    the origin of rivers is on the \emph{other side} of the ranges, not
    \emph{in} the ranges
  \item
    all these rivers flow to nowhere because sand has blocked their
    course
  \end{itemize}
\item
  hills
\item
  gorges and mountains

  \begin{itemize}
  \tightlist
  \item
    especially Davenport, Murchison and MacDonnell Ranges
  \item
    note Mount Sonder is 1,200 m high
  \end{itemize}
\end{itemize}

\emph{Plant species}

\emph{Sida} - half the species are undescribed and the other half are
hard to tell apart

Dominant plant families

\begin{longtable}[t]{lr}
\toprule
Family & No. of species\\
\midrule
Poaceae & 320\\
Asteraceae & 184\\
Fabaceae & 174\\
Chenopodeaceae & 146\\
Malvaceae & 136\\
\bottomrule
\end{longtable}

Absent species

\begin{itemize}
\tightlist
\item
  no Rutaceae species!
\item
  no orchids!
\end{itemize}

\emph{Random bits}

\begin{itemize}
\tightlist
\item
  \emph{Callitris} only lives on slopes, away from fire prone areas
  (e.g.~up on Simpsons Gap)
\item
  Wind pollination is ore common than in other regions (i.e.~more common
  than insect pollination)
\item
  Some plants were here when Central Australia was a lot wetter
  (e.g.~cycads - \emph{Macrozamia} and ferns) - these are \emph{refugia}
  species
\item
  Plant distribution patterns depend on the origins of plants:

  \begin{itemize}
  \tightlist
  \item
    temperate
  \item
    tropical
  \item
    regugia
  \end{itemize}
\end{itemize}

(Note : See paper

\end{document}
